\documentclass[11pt]{article}

\usepackage[margin=1in]{geometry}
\usepackage{amsmath,amssymb,bm,mathtools}
\usepackage{enumitem}
\usepackage{hyperref}

\hypersetup{colorlinks=true,linkcolor=blue,urlcolor=blue}

\newcommand{\mtrip}{\bm{m}}
\newcommand{\etrip}{\bm{e}}
\newcommand{\Mglob}{\bm{M}_{\mathrm{global}}}
\newcommand{\uone}{\bm{u}_1}
\newcommand{\Fmu}{\mathcal{F}_\mu}
\newcommand{\Lie}{\mathcal{L}}
\newcommand{\sgnF}{\sigma_F}
\newcommand{\sgnC}{\sigma_C}
\newcommand{\sgnA}{\sigma_A}
\newcommand{\sgnG}{\sigma_G}

\begin{document}

\title{Symbolic Audit of the Implemented $\chi/W$ Surface for the\texorpdfstring{ $H\parallel c$, $E\parallel a$ }{ H//c, E//a } Readout in TmFeO$_3$}
\author{}
\date{May 2026}

\maketitle

\section{Statement of the audit}

The question is whether the \emph{currently implemented and currently scanned}
mixed Fe--Tm Hamiltonian can generate a physically meaningful Fe--Tm crosspeak
at excitation near the Fe qAFM mode and detection near the Tm
$E_1 \rightarrow E_2$ line, when the experimental readout is taken from the
stored mixed-lattice observable \texttt{M\_global\_SU3} in its
\texttt{lambda1/lambda2} components.

The working criterion is the same as in the previous note: amplitudes of order
$10^{-3}$ or smaller are not counted as genuine.

The conclusion remains negative, but the correct reason is slightly sharper than
the markdown draft: the obstruction is not that one is measuring a naive
uniform sum of raw local $\lambda_1$ and $\lambda_2$. The real obstruction is
that the code measures an \emph{induced-frame} global SU(3) observable, and in
the present Geometry-II setting the implemented $\chi$ source, the CEF
Liouvillian, and the scanned $W_1^{xy}/W_1^{yz}$ vertices never generate the
four-sublattice character required by that readout.

\section{The actual code-level readout operator}

In the mixed Tm sector the active magnetic and electric SU(3) triplets are
\begin{equation}
\mtrip_\mu =
\begin{pmatrix}
\lambda_\mu^2 \\
\lambda_\mu^5 \\
\lambda_\mu^7
\end{pmatrix},
\qquad
\etrip_\mu =
\begin{pmatrix}
\lambda_\mu^1 \\
\lambda_\mu^4 \\
\lambda_\mu^6
\end{pmatrix},
\end{equation}
with one triplet per Tm sublattice $\mu=1,\dots,4$.

The code does \emph{not} store a raw local-site sum. Instead it uses the
induced frame
\begin{equation}
\begin{aligned}
\Fmu &= \mu_{\mathrm{act}}^{-1} R_\mu \, \mu_{\mathrm{act}}, \\
R_1 &= \mathrm{diag}(+,+,+),
&R_2 &= \mathrm{diag}(+,-,-), \\
R_3 &= \mathrm{diag}(-,+,-),
&R_4 &= \mathrm{diag}(-,-,+),
\end{aligned}
\label{eq:induced-frame}
\end{equation}
with the same $3\times 3$ active block acting on both
$\mtrip_\mu$ and $\etrip_\mu$.

The stored global observables are therefore
\begin{equation}
\Mglob^{(m)} = \frac{1}{4}\sum_{\mu=1}^{4} \Fmu \, \mtrip_\mu,
\qquad
\Mglob^{(e)} = \frac{1}{4}\sum_{\mu=1}^{4} \Fmu \, \etrip_\mu,
\label{eq:stored-global}
\end{equation}
and the saved components labelled \texttt{lambda2} and \texttt{lambda1} are
just
\begin{equation}
M_{\mathrm{global}}[\lambda_2] = \uone^{\mathsf T}\Mglob^{(m)},
\qquad
M_{\mathrm{global}}[\lambda_1] = \uone^{\mathsf T}\Mglob^{(e)},
\label{eq:stored-lambda12}
\end{equation}
with $\uone=(1,0,0)^{\mathsf T}$.

This matters because Eqs.~\eqref{eq:stored-global}--\eqref{eq:stored-lambda12}
already include the four-sublattice sign structure. For example, the default
projected moment matrix makes the $c$-axis magnetic probe proportional to the
stored global $\lambda_2$ component, which in local language is the
$(+,-,-,+)$-weighted sum of raw local $\lambda_2$ values, not a naive uniform
sum.

\section{Implemented Hamiltonian in the present scan surface}

At the level relevant here, the mixed model is
\begin{equation}
H = H_{\mathrm{CEF}} + H_\chi + H_W + H_{\mathrm{drive}},
\end{equation}
with
\begin{equation}
H_{\mathrm{CEF}} = \sum_{\mu=1}^{4}\left(\alpha \, \lambda_\mu^3 + \beta \, \lambda_\mu^8\right),
\label{eq:hcef}
\end{equation}
\begin{equation}
H_\chi = \sum_{i,\mu} S_i^a \, \chi_{ab}(i,\mu) \, \lambda_\mu^b,
\qquad
b \in \{2,5,7\},
\label{eq:hchi}
\end{equation}
\begin{equation}
H_W = \sum_{i,\mu} \lambda_\mu^{(n)} \, W_n^{ab}(i,\mu) \, S_i^a S_i^b,
\qquad
n \in \{1,3,4,6,8\}.
\label{eq:hw}
\end{equation}

The code-level support is therefore:
\begin{itemize}[leftmargin=*]
\item $\chi$ populates only the magnetic triplet $(\lambda_2,\lambda_5,\lambda_7)$,
\item $W$ never populates $\lambda_2$ at all,
\item the present scan surface used in the existing 2DCS campaign turns on pure
      $\chi$ lines and the two even $W_1$ lines $W_1^{xy}$ and $W_1^{yz}$.
\end{itemize}

The odd $W_{4,6}$ channels exist in the builder, but they are not part of the
current scan surface audited here.

\section{Klein-four character bookkeeping and the SU(3) EOM}

For the symbolic audit it is convenient to expand any four-sublattice pattern in
the Klein-four character basis
\begin{equation}
\sgnF=(+,+,+,+),\quad
\sgnC=(+,+,-,-),\quad
\sgnA=(+,-,+,-),\quad
\sgnG=(+,-,-,+).
\label{eq:characters}
\end{equation}

For any site-resolved SU(3) component $x_\mu^a$ we write
\begin{equation}
x_\mu^a = \sum_{X\in\{F,C,A,G\}} \sigma_X(\mu) \, x_X^a,
\qquad
x_X^a = \frac{1}{4}\sum_{\mu=1}^{4}\sigma_X(\mu) \, x_\mu^a.
\label{eq:character-decomposition}
\end{equation}

This is only a bookkeeping decomposition of the four-site pattern. It is not a
claim that the full Tm sector should always be reduced to a single Fe-like
Bertaut variable.

With $[\lambda^a,\lambda^b]=2 i f^{abc}\lambda^c$, the local SU(3) equation of
motion takes the form
\begin{equation}
\dot{\ell}_\mu^a = \kappa \, f^{abc} \, \ell_\mu^b h_\mu^c,
\label{eq:local-eom}
\end{equation}
where $\kappa$ is the overall convention-dependent prefactor and
$h_\mu^c = -\partial H / \partial \ell_\mu^c$.

Substituting Eq.~\eqref{eq:character-decomposition} into
Eq.~\eqref{eq:local-eom} gives the character-resolved equation
\begin{equation}
\dot{\ell}_Z^a
=
\kappa \sum_{\substack{X,Y\in\{F,C,A,G\}\\XY=Z}}
f^{abc} \, \ell_X^b h_Y^c,
\label{eq:character-eom}
\end{equation}
where the multiplication is the Klein-four group law coming from pointwise sign
products, for example $G\cdot A=C$ and $G\cdot G=F$.

Equation~\eqref{eq:character-eom} is the key symbolic audit. It implies three
facts immediately:
\begin{enumerate}[leftmargin=*]
\item An $F$-character field preserves the four-sublattice character, because
      $F\cdot X = X$.
\item To convert a $G$-character coherence into a $C$-character coherence one
      needs an $A$-character field, because $G\cdot A = C$.
\item The subset $\{F,G\}$ is closed under multiplication, so a dynamics built
      only from $F$- and $G$-character fields can never generate a
      $C$-character response.
\end{enumerate}

The CEF field in Eq.~\eqref{eq:hcef} is uniform, hence pure $F$-character.
Therefore the linearized CEF Liouvillian is block diagonal in the character
label. In particular, on the $1 \leftrightarrow 2$ coherence pair one has
schematically
\begin{equation}
\partial_t
\begin{pmatrix}
\ell_X^1 \\
\ell_X^2
\end{pmatrix}
=
\begin{pmatrix}
0 & +\Omega_{12} \\
-\Omega_{12} & 0
\end{pmatrix}
\begin{pmatrix}
\ell_X^1 \\
\ell_X^2
\end{pmatrix}
+ \cdots,
\label{eq:12-rotation}
\end{equation}
with the \emph{same} character label $X$ on both components. The CEF can rotate
$\lambda_2$ into $\lambda_1$, but it cannot change the four-sublattice
character while doing so.

\section{Geometry-II source channels}

For the readout geometry of interest, the derivation already recorded in
\texttt{docs/tmfeo3\_notes.tex} gives
\begin{equation}
H_{\mathrm{drive}}^{m,\mathrm{Tm}} \propto h_z \, \Lambda_G^2,
\qquad
H_{\mathrm{drive}}^{e,\mathrm{Tm}} \propto e_x \, \Lambda_C^1.
\label{eq:geometry-two-source}
\end{equation}

Thus the desired Fe$\rightarrow$Tm crosspeak requires the following symbolic
sequence:
\begin{equation}
\begin{aligned}
&\text{qAFM magnetic leg: } G\text{-character in the magnetic triplet} \\
&\hspace{5.9em}\longrightarrow
\text{electric readout leg: } C\text{-character in the electric triplet}.
\end{aligned}
\label{eq:needed-conversion}
\end{equation}

From Eq.~\eqref{eq:character-eom}, this means that some implemented vertex must
provide an \emph{$A$-character} effective field while also supplying the
triplet conversion needed to turn the magnetic $\lambda_2$-type coherence into
the electric $\lambda_1$-type readout direction.

\section{Why the present chi plus scanned W1 surface cannot do this}

\subsection{Chi alone}

The magnetic $\chi$ tensor populates only the magnetic triplet and injects the
qAFM-coupled $\lambda_2$ channel. Equation~\eqref{eq:12-rotation} allows that
source to rotate locally into a $\lambda_1$-type coherence, but only inside the
same character block. Therefore a $G$-character $\chi$ source can only generate
$G$-character $\lambda_1/\lambda_2$ response. It cannot generate the required
$C$-character electric readout of Eq.~\eqref{eq:geometry-two-source}.

So $\chi$ provides the right frequency pair, but the wrong character sector.

\subsection{W1xy and W1yz}

The present scan surface turns on only the even $W_1$ channel. Linearizing its
Fe bilinear around the $\Gamma_2$ background
\begin{equation}
\Gamma_2: \qquad F_x \oplus G_z,
\end{equation}
and the qAFM fluctuation
\begin{equation}
\delta Q_{\mathrm{qAFM}} \sim F_z \oplus G_y,
\end{equation}
gives the character content
\begin{equation}
S_x \delta S_y : F \cdot G = G,
\qquad
S_z \delta S_y : G \cdot G = F.
\label{eq:w1-characters}
\end{equation}

Hence, at the level linear in the qAFM normal coordinate,
\begin{itemize}[leftmargin=*]
\item $W_1^{xy}$ furnishes only a $G$-character electric-triplet source,
\item $W_1^{yz}$ furnishes only an $F$-character electric-triplet source.
\end{itemize}

Neither one supplies the missing $A$-character vertex. Since the CEF field is
also $F$-character, the entire symbolic closure of the present
$\chi + W_1^{xy} + W_1^{yz} + H_{\mathrm{CEF}}$ surface is
\begin{equation}
\{F,G\} \cdot \{F,G\} = \{F,G\},
\end{equation}
which never reaches the $C$ block required by the $E\parallel a$ readout.

This is the cleanest symmetry statement of the obstruction: with the current
scan surface, the EOM never leaves the $\{F,G\}$ subgroup, while the target
readout needs the $C$ block.

\subsection{Why this is enough to rule out the observed crosspeak}

The desired third-order signal is not merely a Tm resonance near
$\omega_{12}$. It is specifically a Geometry-II Fe$\rightarrow$Tm crosspeak:
qAFM excitation on the magnetic leg and electric-triplet readout on the Tm leg.

The current implementation fails to provide the required symbolic bridge:
\begin{enumerate}[leftmargin=*]
\item $\chi$ provides the qAFM-active magnetic source but remains trapped in the
      wrong character block.
\item $W_1$ provides electric-triplet weight but only in $F$ or $G$, never in
      the $C$ block demanded by $E\parallel a$.
\item The uniform CEF field preserves character and therefore cannot repair the
      mismatch.
\end{enumerate}

Therefore the present Hamiltonian does not contain a path from the qAFM-fed
Tm sector into the actual induced-frame Geometry-II readout operator.

\section{Consistency with the existing 2DCS data}

The symbolic audit matches the cached spectra already generated in the repo.
For the baselines and for the full $W_1^{xy}/W_1^{yz}$ grid,
\texttt{M\_NL\_SU3\_FF\_lambda1} and
\texttt{M\_NL\_SU3\_FF\_lambda2} are exactly zero at the nearest grid points to
\begin{equation}
\omega_\tau = 0.911105\;\mathrm{THz},
\qquad
\omega_t = 0.493319\;\mathrm{THz},
\qquad
\omega_t = 0.693729\;\mathrm{THz}.
\end{equation}

The separate raw-state control also remains consistent with the EOM audit:
it can hide tiny weight in a non-measured character channel, but that weight is
only $6.99\times 10^{-5}$ at $(0.911,0.493)$ THz and therefore still far below
the $10^{-3}$ threshold.

\section{Final verdict}

For the current TmFeO$_3$ implementation and the current scan surface, the
absence of the Fe--Tm crosspeak is structural.

The decisive symbolic facts are:
\begin{enumerate}[leftmargin=*]
\item the code-level readout lives in the induced-frame observable
      \texttt{M\_global\_SU3}, not in a raw local-site sum,
\item the Geometry-II source channels require a $G \rightarrow C$ transfer in
      the Tm sector,
\item the present $\chi$ plus scanned $W_1$ surface supplies only $F$- and
      $G$-character vertices, while the CEF Liouvillian preserves character,
\item therefore the linearized SU(3) EOM never reaches the required $C$ block.
\end{enumerate}

So the observed qAFM $\rightarrow$ $E_1 \rightarrow E_2$ Fe--Tm crosspeak
cannot be produced by retuning the existing $\chi$ coefficients or the current
$W_1^{xy}/W_1^{yz}$ scan line alone. To obtain that signal, the model needs an
additional interaction or readout channel that injects the missing character
structure into the electric triplet in the induced-frame basis.

\end{document}